\vspace{-0.1em}
The process of star formation within galaxies is fueled by gas from the interstellar medium (ISM), which exists in molecular, atomic and ionised phases. This study investigates the relationship between star formation and the different phases of the ISM in a sample of 16 nearby galaxies, composed of 9 late-type/dwarfs and 7 spirals.
We used CO data from \textit{HERACLES}, atomic hydrogen (\HI) from \textit{THINGS}, 24 $\mu$m from \textit{SINGS}, and UV from \textit{GALEX NGS} surveys. We decomposed the \HI\ emission spectra using the ROHSA algorithm, which takes into account the spatial coherence of neighbouring spectra, to separate kinematically narrow and broad \HI\ components. We use these components as kinematic proxies for the cold neutral medium (CNM) and warm neutral medium (WNM), rather than direct measurements of the thermal phases.\\
We examine the spatial distribution of the kinematically cold and warm components and find that both components extend across the outermost observed atomic discs. The narrow component contains a median of about 45\% of the total \HI\ mass and has a more filamentary appearance than the broad component. Its fractional contribution is not consistently higher in stellar spiral arms, contrary to theoretical expectations that cooling in the arms could favour cold gas formation. In addition, we find no clear evidence that narrow \HI\ is generally more concentrated near CO emissions than broad \HI, as would be expected if the narrow component were more closely associated with molecular gas.\\
We examine the Kennicutt--Schmidt relation for the different gas components to investigate which components are more closely linked to star formation. In dwarf galaxies, the kinematically cold component shows a stronger overall correlation with the SFR than the total \HI\ or the warm component. In spiral galaxies, the warm component is more globally correlated. A possible explanation of this discrepancy is that line width traces different gas conditions in the two types of galaxies. In dwarfs, narrow emission likely traces cooler, less turbulent gas associated with conditions favourable for star formation. In spirals, star-forming regions can contain cold gas whose lines are broadened by turbulence or unresolved motions. Overall, molecular gas remains the component that is more closely linked to the SFR, as previously shown in the literature. The role of kinematically cold atomic gas in the \HI-to-H$_2$ transition remains ambiguous at our sub-kpc resolution.

\noindent\textbf{Keywords:} Nearby galaxies, Interstellar medium, Star formation
\newpage